\chapter{The rest of the language}
\label{ch:forms}

Five more ideas and you know all of \nodwyn{} Scheme: functions that take
any number of arguments, errors you can catch, templates for building lists,
vectors, and closures that remember.

\section{Any number of arguments}

A parameter list with a dot collects the extra arguments into a list:

\begin{scm}
(define (describe . plants) (length plants))
(describe "a" "b" "c")                        ;=> 3

(define (greet name . rest)
  (if (null? rest) name (string-append name " & " (car rest))))
(greet "basil" "mint")                        ;=> "basil & mint"

((lambda args args) 1 2 3)                    ;=> (1 2 3)
\end{scm}

\nodwyn's own \code{(sub-agent system task tools . name)} uses this for its
optional last argument.

\section{Errors and \texttt{guard}}

\code{(error message irritant \ldots)} raises an error. \code{guard} catches
one: it binds the error to a name and tries clauses like \code{cond}:

\begin{scm}
(define (safe-div a b)
  (guard (e (#t "cannot divide"))
    (/ a b)))
(safe-div 1 0)                                ;=> "cannot divide"
\end{scm}

The error is an association list with a \code{message} and the
\code{irritants}:

\begin{scm}
(guard (e ((string-contains? (assoc-ref e 'message) "boom") 'caught)
          (else (assoc-ref e 'irritants)))
  (error "boom" 1 2))                         ;=> caught

(guard (e ((assoc-ref e 'message) => string-upcase))
  (error "loud"))                             ;=> "LOUD"
\end{scm}

If no clause matches, the error keeps going up, as if the \code{guard} was
not there:

\begin{scm}
(guard (e ((string? e) e))    ; e is an alist, never a string
  (error "fizz"))             ; error: fizz
\end{scm}

\begin{careful}
Two things pass straight through every \code{guard}, whatever its clauses
say: a \textbf{denial} by the guard (a primitive the guard refused) and a
\textbf{cancellation} (you pressed Ctrl-C, or a timeout fired). So do the
loop's hard stops: a spent budget, the tool-call limit, a progress stop.
A program cannot hide that it was refused, and an abort always unwinds.
That is why a strategy with a \code{guard (e (else \ldots))} around
everything still cannot keep a runaway turn alive.
\end{careful}

\section{Quasiquote: templates for lists}

Quote gives you a list exactly as written. \textbf{Quasiquote}, a
backquote, gives you a list \emph{template}: a comma evaluates one piece,
and comma-at splices a list in:

\begin{scm}
(define n 3)
`(n is ,n and doubled ,(* 2 n))          ;=> (n is 3 and doubled 6)
`(plants ,@(list "basil" "mint") end)    ;=> (plants "basil" "mint" end)
\end{scm}

It is perfect for the metadata \nodwyn{} asks for when you register a tool:

\begin{scm}
(define doc "Water a plant.")
`((doc . ,doc) (params . ()))            ;=> ((doc . "Water a plant.") (params))
\end{scm}

(Only one level of quasiquote is supported; you will not miss the second.)

\section{Vectors}

A vector is a fixed-length, indexed collection, and the one thing in the
language you can change in place:

\begin{scm}
(define v (vector "a" "b" "c"))
(vector-set! v 1 "B")
v                                        ;=> #("a" "B" "c")
(vector-length v)                        ;=> 3
(vector->list #(1 2 3))                  ;=> (1 2 3)
(list->vector '(1 2))                    ;=> #(1 2)
(vector-ref #(1 2 3) 5)    ; error: index 5 out of range for vector of length 3
\end{scm}

Reach for a vector when you need to jump to the \emph{n}th element quickly
or fill a table as you go; reach for a list otherwise.

\section{Closures: functions that remember}

A \code{lambda} remembers the variables around it when it was made. Combine
that with \code{set!} and you get private state:

\begin{scm}
(define (make-counter)
  (let ((n 0))
    (lambda ()
      (set! n (+ n 1))
      n)))

(define a (make-counter))
(define b (make-counter))
(a) (a)
(list (a) (b))                           ;=> (3 1)
\end{scm}

\code{a} and \code{b} each have their own \code{n}, and nothing else in the
program can touch it. This is the pattern behind every \nodwyn{} habit that
wraps a loop function: it keeps the original in a variable of its own, and
its new version closes over it (Chapter~\ref{ch:strategies}).

\section{The smaller forms}

\begin{scm}
(case 'tomato
  ((basil mint) "herb")
  ((tomato bean) "vegetable")
  (else "unknown"))                      ;=> "vegetable"

(when (> 2 1) "ran")                     ;=> "ran"
(unless (> 2 1) "ran")                   ; nothing: returns void

(begin (define a 1) (define b 2) (+ a b));=> 3

(define (hyp a b)
  (define (sq x) (* x x))                ; a definition inside a body
  (sqrt (+ (sq a) (sq b))))
(hyp 3 4)                                ;=> 5.0

(letrec ((even? (lambda (n) (if (= n 0) #t (odd? (- n 1)))))
         (odd?  (lambda (n) (if (= n 0) #f (even? (- n 1))))))
  (even? 100001))                        ;=> #f
\end{scm}

\code{case} compares with \code{equal?}, so strings work as keys too.
\code{letrec} lets its functions call each other; notice that
\code{even?} and \code{odd?} call each other in tail position, so the
hundred thousand hops take no stack at all.

\section{Sameness}

\begin{center}
\begin{tabularx}{\linewidth}{lX}
\toprule
\code{equal?} & same shape and contents. Use this one. \\
\code{eq?}, \code{eqv?} & the same object, for pairs, vectors and procedures; the same value for everything else, including strings \\
\bottomrule
\end{tabularx}
\end{center}

\begin{scm}
(equal? '(1 (2 3)) (list 1 (list 2 3)))  ;=> #t
(eq? (list 1) (list 1))                  ;=> #f   two different pairs
(eq? 'a 'a)                              ;=> #t
(equal? #(1 2) #(1 2))                   ;=> #t
\end{scm}

\section{What is deliberately missing}

No \code{call/cc}, no macros (\code{define-syntax}), no \code{eval}, no
characters, no string mutation, no \code{set-car!}, no fractions or
bignums, no nested quasiquote. The language is small on purpose: it is
there to express policy, and a small language is one you can hold in your
head and one whose every program the guard can reason about. There is
also, deliberately, no way for a Scheme program to reach the kernel's
evaluator: loading code can add functions and primitives, never new
special forms.

\begin{puzzle}[ 8: a safe number parser]
Write \code{(parse-ml text)} that turns \code{"250 ml"} into \code{250},
and returns \code{\#f} for text that does not start with a number. Then
write \code{(must-parse-ml text)} that raises an error naming the text
when parsing fails, and show how a caller catches it with
\code{guard}.\ans{8}
\end{puzzle}

\begin{puzzle}[ 9: a watering log]
Write \code{(make-log)} that returns a function. Called with a string, it
remembers it; called with no arguments, it returns everything remembered,
oldest first. Use a closure and \code{set!}.\ans{9}
\end{puzzle}
